% =============================================================================
% Appendix A --- Full FSM State Diagrams
% Derived from Context 2 §6 Mermaid diagrams. Figure placeholders with
% explicit transition-condition captions so the diagrams can be drawn
% without ambiguity.
% =============================================================================

\chapter{Full FSM State Diagrams}
\label{app:fsm-diagrams}

This appendix collects full-page, fully-labeled renderings of the four
state machines described in \Cref{ch:rtl}. Each figure is a placeholder:
when the corresponding diagram file is dropped into \sv{figures/}, the
chapter will compile unchanged. The transition conditions listed in each
caption are the authoritative conditions taken from \sv{controller.sv}
and should be reproduced verbatim on the diagrams.

\section{Main FSM}
\label{app:fsm-main}

\begin{figure}[htbp]
    \centering
    \includegraphics[width=0.95\linewidth]{fig_app_main_fsm_full}
    \caption[Main FSM of the ANN controller (fully labeled)]
        {Full main FSM of \modname{ann\_controller}. Transition
         conditions (verbatim from \sv{controller.sv}):
         \stateval{S\_IDLE}\,\(\rightarrow\)\,\stateval{S\_PROGRAM} on
         \signame{valid}~\sv{\&\&}~(\signame{cmd}~=~\sv{CMD\_PROG});
         \stateval{S\_IDLE}\,\(\rightarrow\)\,\stateval{S\_ERASE} on
         \signame{valid}~\sv{\&\&}~(\signame{cmd}~=~\sv{CMD\_ERASE});
         \stateval{S\_IDLE}\,\(\rightarrow\)\,\stateval{S\_READ} on
         \signame{valid}~\sv{\&\&}~(\signame{cmd}~=~\sv{CMD\_READ});
         \stateval{S\_IDLE}\,\(\rightarrow\)\,\stateval{S\_COLLECT\_DATA}
         on \signame{valid}~\sv{\&\&}~(\signame{cmd}~=~\sv{CMD\_INF});
         self-loop otherwise (\(!\)\signame{valid} or default
         command path);
         \stateval{S\_RESET}\,\(\rightarrow\)\,\stateval{S\_PROGRAM}
         unconditionally (note: no edge leads into
         \stateval{S\_RESET} in the current RTL);
         \stateval{S\_PROGRAM}\,\(\rightarrow\)\,\stateval{S\_VERIFY}
         when \signame{prog\_state}~=~\sv{PROG\_COMPLETE};
         \stateval{S\_VERIFY}\,\(\rightarrow\)\,\stateval{S\_IDLE}
         when \signame{verify\_state}~=~\sv{VERIFY\_DONE};
         \stateval{S\_VERIFY}\,\(\rightarrow\)\,\stateval{S\_PROGRAM}
         on under-programmed readback with
         \signame{prog\_retry\_cnt}~\(<\)~\sv{MAX\_PROG\_RETRIES};
         \stateval{S\_VERIFY}\,\(\rightarrow\)\,\stateval{S\_ERASE}
         on over-programmed readback or exhausted retries;
         \stateval{S\_ERASE}\,\(\rightarrow\)\,\stateval{S\_IDLE}
         on \sv{ERASE\_COMPLETE} with
         (\signame{erase\_from\_host}~\(||\)~\(!\)(\signame{retry\_cnt}~\(<\)~3));
         \stateval{S\_ERASE}\,\(\rightarrow\)\,\stateval{S\_PROGRAM}
         on \sv{ERASE\_COMPLETE} with
         \(!\)\signame{erase\_from\_host}~\sv{\&\&}~(\signame{retry\_cnt}~\(<\)~3);
         \stateval{S\_READ}\,\(\rightarrow\)\,\stateval{S\_IDLE} on
         \signame{pulse\_done}, else self-loop;
         \stateval{S\_COLLECT\_DATA}\,\(\rightarrow\)\,\stateval{S\_COLLECT\_DATA}
         on \signame{valid}~\sv{\&\&}~(\signame{data\_count}~\(<\)~7) or
         when \signame{valid} is low;
         \stateval{S\_COLLECT\_DATA}\,\(\rightarrow\)\,\stateval{S\_COMPUTE}
         on \signame{valid}~\sv{\&\&}~\(!\)(\signame{data\_count}~\(<\)~7);
         \stateval{S\_COMPUTE}\,\(\rightarrow\)\,\stateval{S\_RESULT}
         on \signame{pulse\_done}, else self-loop;
         \stateval{S\_RESULT}\,\(\rightarrow\)\,\stateval{S\_IDLE}
         unconditionally.}
    \label{fig:app-main-fsm-full}
\end{figure}

\clearpage

\section{Program Sub-FSM}
\label{app:fsm-prog}

\begin{figure}[htbp]
    \centering
    \includegraphics[width=0.95\linewidth]{fig_app_prog_subfsm_full}
    \caption[Program sub-FSM (fully labeled)]{Full Program sub-FSM
        inside \stateval{S\_PROGRAM}. Reset: \sv{PROG\_HIZ} on
        \signame{!rst\_n}. Transitions:
        \sv{PROG\_HIZ}\,\(\rightarrow\)\,\sv{PROG\_SELECT} unconditionally;
        \sv{PROG\_SELECT}\,\(\rightarrow\)\,\sv{PROG\_WRITE} on
        \signame{op\_done}, else self-loop;
        \sv{PROG\_WRITE}\,\(\rightarrow\)\,\sv{PROG\_WAIT\_ACK} on
        \signame{pulse\_done}, else self-loop;
        \sv{PROG\_WAIT\_ACK}\,\(\rightarrow\)\,\sv{PROG\_COMPLETE} on
        \signame{op\_done}, else self-loop;
        \sv{PROG\_COMPLETE} issues \signame{next\_state}~=~\stateval{S\_VERIFY}
        and \signame{next\_prog\_state}~=~\sv{PROG\_HIZ} simultaneously,
        returning to the main FSM which enters \stateval{S\_VERIFY} on
        the next clock edge.}
    \label{fig:app-prog-subfsm-full}
\end{figure}

\clearpage

\section{Verify Sub-FSM}
\label{app:fsm-verify}

\begin{figure}[htbp]
    \centering
    \includegraphics[width=0.95\linewidth]{fig_app_verify_subfsm_full}
    \caption[Verify sub-FSM (fully labeled)]{Full Verify sub-FSM inside
        \stateval{S\_VERIFY}. Reset: \sv{VERIFY\_IDLE}. Transitions
        (\signame{verify\_next} is the combinational output that the
        sequential block samples each cycle):
        \sv{VERIFY\_IDLE}\,\(\rightarrow\)\,\sv{VERIFY\_CHECK} when
        \sv{PULSE\_TOTAL\_READ}~\(\le\)~\sv{8'd1};
        \sv{VERIFY\_IDLE}\,\(\rightarrow\)\,\sv{VERIFY\_WAIT} otherwise;
        \sv{VERIFY\_WAIT}\,\(\rightarrow\)\,\sv{VERIFY\_CHECK} when
        \signame{verify\_pulse\_idx}~=~(\sv{PULSE\_TOTAL\_READ[7:0]~-~8'd1}),
        else self-loop;
        \sv{VERIFY\_CHECK}\,\(\rightarrow\)\,\sv{VERIFY\_DONE} when
        \signame{weight\_read\_data}~=~\signame{expected\_weight};
        \sv{VERIFY\_CHECK}\,\(\rightarrow\)\,\sv{VERIFY\_IDLE} with
        \signame{next\_state}~=~\stateval{S\_PROGRAM},
        \signame{next\_prog\_state}~=~\sv{PROG\_HIZ} when
        \signame{weight\_read\_data}~\(<\)~\signame{expected\_weight}
        and \signame{prog\_retry\_cnt}~\(<\)~\sv{MAX\_PROG\_RETRIES};
        \sv{VERIFY\_CHECK}\,\(\rightarrow\)\,\sv{VERIFY\_IDLE} with
        \signame{next\_state}~=~\stateval{S\_ERASE} when
        \signame{weight\_read\_data}~\(<\)~\signame{expected\_weight}
        and \signame{prog\_retry\_cnt}~\(\ge\)~\sv{MAX\_PROG\_RETRIES};
        \sv{VERIFY\_CHECK}\,\(\rightarrow\)\,\sv{VERIFY\_IDLE} with
        \signame{next\_state}~=~\stateval{S\_ERASE} when
        \signame{weight\_read\_data}~\(>\)~\signame{expected\_weight};
        \sv{VERIFY\_DONE}\,\(\rightarrow\)\,\sv{VERIFY\_IDLE} with
        \signame{next\_state}~=~\stateval{S\_IDLE}.}
    \label{fig:app-verify-subfsm-full}
\end{figure}

\clearpage

\section{Erase Sub-FSM}
\label{app:fsm-erase}

\begin{figure}[htbp]
    \centering
    \includegraphics[width=0.95\linewidth]{fig_app_erase_subfsm_full}
    \caption[Erase sub-FSM (fully labeled)]{Full Erase sub-FSM inside
        \stateval{S\_ERASE}. Transitions:
        \sv{ERASE\_HIZ}\,\(\rightarrow\)\,\sv{ERASE\_SELECT}
        unconditionally;
        \sv{ERASE\_SELECT}\,\(\rightarrow\)\,\sv{ERASE\_PULSE} on
        \signame{op\_done}, else self-loop;
        \sv{ERASE\_PULSE}\,\(\rightarrow\)\,\sv{ERASE\_WAIT\_ACK} on
        \signame{pulse\_done}, else self-loop;
        \sv{ERASE\_WAIT\_ACK}\,\(\rightarrow\)\,\sv{ERASE\_COMPLETE} on
        \signame{op\_done}, else self-loop;
        \sv{ERASE\_COMPLETE} exits with
        \signame{erase\_next}~=~\sv{ERASE\_HIZ} and selects
        \signame{next\_state} according to
        \Cref{tab:erase-complete}: \stateval{S\_IDLE} if
        \signame{erase\_from\_host}; \stateval{S\_PROGRAM} if
        \(!\)\signame{erase\_from\_host}~\sv{\&\&}~(\signame{retry\_cnt}~\(<\)~3);
        \stateval{S\_IDLE} with
        \signame{erase\_max\_retries\_exhausted}~=~1 otherwise.}
    \label{fig:app-erase-subfsm-full}
\end{figure}
